\chapter{时序与存储设计}\label{ch:timing}

\section{时序设计规则}

为使 4096 条 VL 的查表通路和逐字节数据通路都能在 8~ns 内完成，全部 RTL 按表~\ref{tab:rules} 的十条规则组织寄存器、RAM 访问和组合运算。

\begin{table}\centering\small
\caption{时序设计规则}\label{tab:rules}
\begin{tabularx}{\linewidth}{>{\centering\arraybackslash}p{0.8cm}L}
\toprule
编号 & 规则 \\
\midrule
1 & 全设计使用单一时钟。 \\
2 & 逐 VL 数据只存放于 RAM，不得出现按 VL 数展开的触发器阵列或并行比较。 \\
3 & 所有 RAM 仅通过 \code{afdx_ram_sdp} 例化，地址、写数据和使能均直接来自寄存器。小寄存器组须注明用途（载荷 FIFO 8 项、归还 FIFO 4 项、接收描述 32 项）。 \\
4 & RAM 读出先经寄存器，再做加、减、比较；中间最多允许一级不超过 4 选 1 的选择。 \\
5 & 每拍最多一级 32 位算术；不使用乘法器和除法器，地址以位拼接生成。 \\
6 & 多于 8 选 1 的多路选择分两拍完成。 \\
7 & 数据逐拍连续推进；RAM 与算术采用寄存流水，帧头期间执行预读。 \\
8 & 同一 RAM 的同拍同址读写必须互斥。 \\
9 & 模块输出以寄存器为主；组合输出仅限 ready 握手、\code{cfg_ok}，以及 \code{afdx_svc} 的数据选择与接收分发直通，后者只做选择和门控。 \\
10 & 以单一 125~MHz、8~ns 为时钟约束，内部设计目标为综合 setup 裕量不小于 2.0~ns；不设置 false path 或多周期例外。 \\
\bottomrule
\end{tabularx}
\end{table}

规则 2 和 4 是“表加单主状态机”结构的基础：逐 VL 状态全部在 RAM 中，每张表只由一个状态机按固定流水访问，每拍的逻辑深度与 VL 数量无关。规则 5 使 RTC 的时间戳换算和偏差计算也拆成每拍一级加减，服务模块与收发核心均在寄存器之间限制组合运算的级数。

\section{存储预算}

全部存储都由 \code{afdx_ram_sdp} 例化，表~\ref{tab:ram} 按 RTL 数组声明的深度和位宽计算逻辑存储位数；器件存储块的映射用量见测试报告。

\begin{table}\centering\small
\caption{存储预算（逻辑位数）}\label{tab:ram}
\begin{tabularx}{\linewidth}{>{\raggedright\arraybackslash}p{3.6cm}>{\raggedleft\arraybackslash}p{2.8cm}>{\raggedleft\arraybackslash}p{2.2cm}L}
\toprule
表 & 深度$\times$位宽 & 位数 & 说明 \\
\midrule
\code{tx_payload} & 131072$\times$8 & 1048576 & 256 页 $\times$ 512 字节 \\
\code{tx_cfg} & 1024$\times$32 & 32768 & 128 条 VL $\times$ 8 字 \\
\code{tx_desc} & 256$\times$44 & 11264 & 每页一项 \\
\code{tx_state} & 128$\times$63 & 8064 & 每条 VL 一项 \\
\code{tx_mnext/pnext/free} & 3$\times$256$\times$8 & 6144 & 链表与空闲页 FIFO \\
\midrule
\code{rx_map} & 65536$\times$13 & 851968 & 任意 VLID 直接映射 \\
\code{rx_state} & 4096$\times$44 & 180224 & 每条 VL 一项 \\
\code{rx_cfg} & 4096$\times$32 & 131072 & 每条 VL 一字 \\
\code{rx_ring_a/b} & 2$\times$4096$\times$8 & 65536 & 每网 4~KiB \\
\midrule
SNMP 请求、回应 & 2$\times$256$\times$8 & 4096 & \\
615A 五块 & --- & 102400 & 收发报文各 1~KiB，LUR 512~B，LUH 2~KiB，加载区 8~KiB \\
\midrule
合计 & --- & 2442112 & 发送 1106816，接收 1228800，服务 106496 \\
\bottomrule
\end{tabularx}
\end{table}

\code{tx_payload} 存放发送载荷，\code{rx_map} 完成 16 位 VLID 到接收索引的映射。这两张深表均由块 RAM 实现。


\section{流水划分}

收发通路把查表、判决和数据传送分开：查表地址由寄存器驱动，RAM 读出先寄存，下一拍再做算术或比较；组帧与接收解析沿字节流逐拍推进。各模块的寄存边界见表~\ref{tab:pipeline-rule}。

\begin{table}\centering\small
\caption{主要通路的流水划分}\label{tab:pipeline-rule}
\begin{tabularx}{\linewidth}{>{\raggedright\arraybackslash}p{2.6cm}LL}
\toprule
通路 & 组合处理 & 寄存边界与推进方式 \\
\midrule
发送调度 & 到期比较、链表指针更新、候选选择 & 配置和状态读出先寄存；入队、授权和扫描按操作阶段推进（第~\ref{sec:txeng}~节） \\
组帧与 MAC & 帧头选择、校验和、CRC 更新 & 帧头期间预读载荷；每拍输出一个字节，双网共用发送状态机（第~\ref{sec:txbld}、\ref{sec:mactx}~节） \\
接收映射与判决 & VLID 查表、配置检查、序号与冗余窗口比较 & 映射表读出寄存后进入判决，各阶段逐拍完成，结果用于提交或回退（第~\ref{sec:rxtables}~节、表~\ref{tab:rxphase}） \\
服务协议处理 & 字节解析、字段比较、时间戳换算 & 同步读出字节后再解析；时间换算和偏差计算拆成逐拍加减（第~\ref{ch:svc}~章） \\
\bottomrule
\end{tabularx}
\end{table}

\code{cfg_ok} 将配置写与引擎访问错开，避免同拍读写同一表项；具体许可和写入时刻见第~\ref{sec:cfgok}~节。综合资源、建立与保持时间检查结果集中列于测试报告。
